%------------------------------------------------------------------------------%
\documentclass[fleqn,utf8,aspectratio=169,12pt]{beamer}

\usepackage{calculo_varias_variaveis-1}

\title{Curvas Paramétricas}

%------------------------------------------------------------------------------%
\begin{document}

\addtitle

%------------------------------------------------------------------------------%
\section{Introdução}

\framedgraphic{Trajetória}{trajeto_maps}{}

%------------------------------------------------------------------------------%
\begin{frame}{Curvas Paramétricas}
  Função de $\R$ em $\R^n$
  \[
    \pause
    \gamma(t)
    = (x, y)
    \pause
    = \big( f(t), g(t) \big)
  \]
  \pause
  ou
  \[
    \pause
    \gamma(t) =
    \left[\begin{array}{>{\,}c<{\,}}
      x \\[1mm]
      y
    \end{array}\right]
    \pause
    =
    \left[\begin{array}{>{\,}c<{\,}}
      f(t) \\[1mm]
      g(t)
    \end{array}\right]
  \]
  \pause
  ou
  \[
    \gamma(t)
    = x\veci + y\vecj
    \pause
    = f(t)\veci + g(t)\vecj
  \]

  \pause
  Podemos pensar na variável como o tempo e a curva como uma trajetória
\end{frame}

%------------------------------------------------------------------------------%
\section{Curvas Paramétricas}

%------------------------------------------------------------------------------%
\begin{frame}{Definição em 2D}
  Se $x$ e $y$
  \pause
  são dados como funções
  \[
    \pause
    x = f(t)
    \qquad
    y = g(t)
  \]
  \pause
  sobre um intervalo $I$ de valores de $t$,
  \pause
  então o conjunto de pontos
  \[
    (x,\; y) = \big( f(t),\; g(t)\big)
  \]
  \pause
  forma uma \emph{curva paramétrica} em $\R^2$
\end{frame}

%------------------------------------------------------------------------------%
\begin{frame}{Definição em 3D}
  Se $x$, $y$ e $z$ são dados como funções
  \[
    x = f(t)
    \qquad
    y = g(t)
    \qquad
    z = h(t)
  \]
  sobre um intervalo $I$ de valores de $t$,
  então o conjunto de pontos
  \[
    (x,\; y,\; z) = \big( f(t),\; g(t),\; h(t) \big)
  \]
  forma uma \emph{curva paramétrica} em $\R^3$
\end{frame}

%------------------------------------------------------------------------------%
\begin{frame}{Definição Geral}
  Se $x_i$, $i=1, \dots, n$ são as coordenadas no espaço $\R^n$
  \[
    x_i = f_i(t)
  \]
  sobre um intervalo $I$ de valores de $t$,
  então o conjunto de pontos
  \[
    (x_1,\; x_2,\; \dots,\; x_n) = \big( f_1(t),\; f_2(t),\; \dots,\; f_n(t)\big)
  \]
  forma uma \emph{curva paramétrica} em $\R^n$
\end{frame}

%------------------------------------------------------------------------------%
\begin{frame}{Nomenclatura}
  As equações são as \emph{equações paramétricas} da curva

  \pause
  \medskip
  \emph{$t$} \quad é o parâmetro da curva

  \pause
  \emph{$I$} \quad intervalo do parâmetro

  \pause
  \medskip
  Se $I$ é um intervalo fechado \emph{\(a\leq t \leq b\)}
  \pause
  \begin{description}[<+->]
    \item[\(\big(f(a), g(a)\big)\)] é o ponto inicial
    \item[\(\big(f(b), g(b)\big)\)] é o ponto final
  \end{description}
\end{frame}

%------------------------------------------------------------------------------%
\begin{frame}{Notações}
  \[
    (x, y)
    = \big(x(t), y(t)\big)
    = \big(f(t), g(t)\big)
  \]

  \pause
  \[
      \begin{pmatrix}  x    \\ y    \end{pmatrix}
    = \begin{pmatrix}  x(t) \\ y(t) \end{pmatrix}
    = \begin{pmatrix}  f(t) \\ g(t) \end{pmatrix}
  \]

  \pause
  \[
    (x_1, x_2, \dots, x_n) = \big( f_1(t), f_2(t), \dots, f_n(t)\big)
  \]
\end{frame}

%------------------------------------------------------------------------------%
\section{Exemplos}

%------------------------------------------------------------------------------%
\begin{question}
  Trace a curva definida pelas equações paramétricas
  \[
    x = t^2       \qquad
    y = t + 1     \qquad
    -\infty < t < \infty
  \]
\end{question}

%------------------------------------------------------------------------------%
\begin{resolution}{Solução}
\begin{columns}
\column{0.4\linewidth}
  Note que

  \pause
  \bigskip
  $y$ varia com $t$ e

  \pause
  $x$ varia com $t^2$

  \pause
  \bigskip
  então a curva é uma parábola

\column{0.5\linewidth}
  \pause
  Calculando alguns pontos
  \[
    \begin{array}{r<{\qquad}r<{\quad}r}
      \hline
      t  & x & y  \\ \hline
      -3 & 9 & -2 \\
      -2 & 4 & -1 \\
      -1 & 1 & 0  \\
      0  & 0 & 1  \\
      1  & 1 & 2  \\
      2  & 4 & 3  \\
      3  & 9 & 4  \\ \hline
    \end{array}
  \]
\end{columns}
\end{resolution}

%------------------------------------------------------------------------------%
\begin{resolution}{Solução}
  \centering
  \includegraphics{B_02_exemplo_curvas_1}
\end{resolution}

%------------------------------------------------------------------------------%

%------------------------------------------------------------------------------%
%------------------------------------------------------------------------------%
\begin{question}
  Identifique geometricamente a curva do Exemplo 1 eliminando o parâmetro $t$
  \[
    x = t^2        \qquad
    y = t + 1      \qquad
    -\infty < t < \infty
  \]

  \pause
  Transformar a curva paramétrica em $t$ em uma equação envolvendo $x$ e $y$
\end{question}

%------------------------------------------------------------------------------%
\begin{resolution}{Solução}
\begin{columns}
\column{0.3\linewidth}
  \onslide<+->{Eliminar $t$ em
  \begin{align*}
    x & = t^2   \\
    y & = t + 1
  \end{align*}}

\column{0.3\linewidth}
  \begin{align*}
    \onslide<+->{y & = t + 1 \\}
    \onslide<+->{t & = y - 1}
  \end{align*}

\column{0.3\linewidth}
  \begin{align*}
    \onslide<+->{x & = t^2} \\
    \onslide<+->{  & = (y - 1)^2} \\
    \onslide<+->{  & = y^2 - 2y + 1}
  \end{align*}
\end{columns}
\end{resolution}

%------------------------------------------------------------------------------%

%------------------------------------------------------------------------------%
\begin{question}
  Determine se os pontos \((2, 3)\) e \((4, -1)\) pertencem à curva paramétrica
  \[
    x = t^2       \qquad
    y = t + 1     \qquad
    -\infty < t < \infty
  \]
\end{question}

%------------------------------------------------------------------------------%
\begin{resolution}{Solução}
\begin{columns}[t]
\column{0.4\linewidth}
  \onslide<+->{Testando o ponto \((2, 3)\)}
  \begin{align*}
    \onslide<+->{y & = 3} \\[2mm]
    \onslide<+->{t & + 1 = 3}         \\[2mm]
    \onslide<+->{t & = 3 -1}
    \onslide<+->{ = 2}      \\[2mm]
    \onslide<+->{x & = t^2}
    \onslide<+->{ = 2^2}
    \onslide<+->{ = 4}
    \onslide<+->{\neq 2}
  \end{align*}
  \onslide<+->{O ponto não pertence a curva}

\column{0.4\linewidth}
\onslide<+->{Testando o ponto \((4, -1)\)}
\begin{align*}
  \onslide<+->{y & = -1}               \\[2mm]
  \onslide<+->{t & + 1 = -1}           \\[2mm]
  \onslide<+->{t & = -1 -1}
  \onslide<+->{ = -2}       \\[2mm]
  \onslide<+->{x & = t^2}
  \onslide<+->{ = (-2)^2}
  \onslide<+->{ = 4}
\end{align*}
\onslide<+->{O ponto pertence a curva}
\end{columns}
\end{resolution}

%------------------------------------------------------------------------------%

%------------------------------------------------------------------------------%
\begin{question}
  Represente graficamente a curva paramétrica
  \[
    x = \cos(t)         \qquad
    y = \sin(t)         \qquad
    0 \leq t \leq 2\pi
  \]
\end{question}

%------------------------------------------------------------------------------%
\begin{resolution}{Solução}
  Sabemos que essa é a parametrização do circulo trigonométrico

  \pause
  Calculando alguns pontos
  \begin{ceqn}
  \[
  \renewcommand*{\arraystretch}{1.5}
  \begin{array}{c<{\quad}r<{\quad}r}
    \hline
    t               & x  & y  \\ \hline
    0               & 1  & 0  \\
    \dfrac{\pi}{2}  & 0  & 1  \\
    \pi             & -1 & 0  \\
    \dfrac{3\pi}{2} & 0  & -1 \\
    2\pi            & 1  & 0  \\ \hline
  \end{array}
  \]
  \end{ceqn}
\end{resolution}

%------------------------------------------------------------------------------%
\framedgraphic*[resolution]{Solução}{B_02_exemplo_4_curva_1}
\framedgraphic*[resolution]{Solução}{B_02_exemplo_4_curva_2}
\framedgraphic*[resolution]{Solução}{B_02_exemplo_4_curva_3}
\framedgraphic*[resolution]{Solução}{B_02_exemplo_4_curva_4}
\framedgraphic*[resolution]{Solução}{B_02_exemplo_4_curva_5}
\framedgraphic*[resolution]{Solução}{B_02_exemplo_4_curva_6}
\framedgraphic*[resolution]{Solução}{B_02_exemplo_4_curva_7}
\framedgraphic*[resolution]{Solução}{B_02_exemplo_4_curva_8}

%------------------------------------------------------------------------------%

%------------------------------------------------------------------------------%
\begin{question}
  Esboce a trajetória definida pela curva paramétrica
  \[
  \begin{cases}
    x = t \cos(t\pi) \\
    y = t \sin(t\pi)
  \end{cases}
  \]
  \(0\leq t\)
\end{question}

%------------------------------------------------------------------------------%
\begin{resolution}{Solução}
  Similar a parametrização do circulo trigonométrico

  \pause
  \begin{ceqn}
  \[
  \renewcommand*{\arraystretch}{1.7}
  \begin{array}{c<{\quad}c<{\quad}c<{\quad}c<{\quad}c<{\quad}c}
    \hline
    t            & t\pi            & \cos(t\pi) & \sin(t\pi) & x  & y             \\ \hline
    0            & 0               & 1          & 0          & 0  & 0             \\
    \dfrac{1}{2} & \dfrac{\pi}{2}  & 0          & 1          & 0  & \dfrac{1}{2}  \\
    1            & \pi             & -1         & 0          & -1 & 0             \\
    \dfrac{3}{2} & \dfrac{3\pi}{2} & 0          & -1         & 0  & -\dfrac{3}{2} \\
    2            & 2\pi            & 1          & 0          & 2  & 0             \\ \hline
  \end{array}
  \]
  \end{ceqn}
\end{resolution}

%------------------------------------------------------------------------------%
\framedgraphic*[resolution]{Solução}{B_02_exemplo_5_curva_1}
\framedgraphic*[resolution]{Solução}{B_02_exemplo_5_curva_2}
\framedgraphic*[resolution]{Solução}{B_02_exemplo_5_curva_3}
\framedgraphic*[resolution]{Solução}{B_02_exemplo_5_curva_4}
\framedgraphic*[resolution]{Solução}{B_02_exemplo_5_curva_5}
\framedgraphic*[resolution]{Solução}{B_02_exemplo_5_curva_6}

%------------------------------------------------------------------------------%

%------------------------------------------------------------------------------%
\begin{question}
  A posição de uma partícula se movendo no plano $xy$ é dada por
  \[
    x = \sqrt{t}  \qquad
    y = t         \qquad
    t \geq 0
  \]

  Esboce sua trajetória
\end{question}

%------------------------------------------------------------------------------%
\begin{resolution}{Solução}
\begin{columns}
\column{0.4\linewidth}
  Eliminando $t$ em
  \begin{align*}
    x & = \sqrt{t} \\
    y & = t
  \end{align*}
  temos
  \[
    x
    = \sqrt{t}
    = \sqrt{y}
  \]

\column{0.4\linewidth}
  Elevando os dois lados ao quadrado
  \begin{align*}
    x   & = \sqrt{y}                \\
    x^2 & = \left(\sqrt{y}\right)^2
          = \abs{y}
          = y                       \\
    y   & = x^2
  \end{align*}
\end{columns}
\end{resolution}

%------------------------------------------------------------------------------%
\begin{resolution}{Solução}
  \centering
  \includegraphics{B_02_exemplo_curvas_6}
\end{resolution}

%------------------------------------------------------------------------------%

%------------------------------------------------------------------------------%
\begin{question}
  Construa uma curva paramétrica cuja trajetória seja uma circunferência\\
  de raio 2 centrada em $(3, 4)$
\end{question}

%------------------------------------------------------------------------------%
\begin{resolution}{Solução}
\begin{columns}[t]
\column{0.3\linewidth}
  Sabemos que
  \begin{align*}
    x & = \cos(t) \\
    y & = \sin(t) \\
    0 & \leq t < 2\pi
  \end{align*}
  parametriza uma circunferência de raio 1 centrada em \((0, 0)\)
\column{0.3\linewidth}
  \pause
  Para que o raio seja 2 multiplicamos as funções por 2
  \begin{align*}
    x & = 2\cos(t) \\
    y & = 2\sin(t) \\
    0 & \leq t < 2\pi
  \end{align*}
\column{0.3\linewidth}
  \pause
  Para mover o centro para $(3, 4)$ somamos esses valores
  \begin{align*}
    x & = 2\cos(t) + 3\\
    y & = 2\sin(t) + 4\\
    0 & \leq t < 2\pi
  \end{align*}
\end{columns}
\end{resolution}

%------------------------------------------------------------------------------%

%------------------------------------------------------------------------------%
\begin{question}
Encontre uma parametrização para o círculo definido pela equação
\[
  (x-3)^2 + (y+1)^2 = 4
\]
\end{question}

%------------------------------------------------------------------------------%
\begin{resolution}{Solução}
  Círculo de raio 2 centrado no ponto \((3, -1)\)

  Parametrização do círculo unitário centrado na origem
  \begin{align*}
    x      & = \cos(\theta) \\
    y      & = \sin(\theta) \\
    \theta & \in [0, 2\pi)
  \end{align*}
\end{resolution}

%------------------------------------------------------------------------------%
\begin{resolution}{Solução}
  Parametrização do círculo de raio 2 centrado na origem
  \begin{align*}
    x      & = 2\cos(\theta) \\
    y      & = 2\sin(\theta) \\
    \theta & \in [0, 2\pi)
  \end{align*}
\end{resolution}

%------------------------------------------------------------------------------%
\begin{resolution}{Solução}
  Parametrização do círculo de raio 2 centrado no ponto \((3, -1)\)
  \begin{align*}
  x      & = 2\cos(\theta) + 3 \\
  y      & = 2\sin(\theta) - 1 \\
  \theta & \in [0, 2\pi)
  \end{align*}
\end{resolution}

%------------------------------------------------------------------------------%

%------------------------------------------------------------------------------%
\begin{question}
Encontre uma parametrização para o movimento de uma partícula que começa no
ponto \((-2, 0)\) e traça a metade superior do círculo \(x^2+y^2=4\)
\end{question}

%------------------------------------------------------------------------------%
\begin{resolution}{Solução}
  Metade superior do círculo de raio 2 centrado na origem

  Parametrização do círculo unitário centrado na origem
  \begin{align*}
    x      & = \cos(\theta) \\
    y      & = \sin(\theta) \\
    \theta & \in [0, 2\pi)
  \end{align*}
\end{resolution}

%------------------------------------------------------------------------------%
\begin{resolution}{Solução}
  Parametrização do círculo de raio 2 centrado na origem
  \begin{align*}
    x      & = 2\cos(\theta) \\
    y      & = 2\sin(\theta) \\
    \theta & \in [0, 2\pi)
  \end{align*}
\end{resolution}

%------------------------------------------------------------------------------%
\begin{resolution}{Solução}
  Parametrização da metade superior
  \begin{align*}
    x      & = 2\cos(\theta) \\
    y      & = 2\sin(\theta) \\
    \theta & \in [0, \pi]
  \end{align*}
\end{resolution}

%------------------------------------------------------------------------------%
\begin{resolution}{Solução}
  Essa parametrização começa em \((2, 0)\) e termina em \((-2, 0)\)


  Queremos inverter o sentido em que percorremos a trajetória
  \begin{align*}
    x      & = 2\cos(\pi -\theta) =    - 2 \cos(\theta)\\
    y      & = 2\sin(\pi -\theta) = \hpm 2 \sin(\theta)\\
    \theta & \in [0, \pi)
  \end{align*}
\end{resolution}

%------------------------------------------------------------------------------%


%------------------------------------------------------------------------------%
\section{Lista Mínima}

%------------------------------------------------------------------------------%
\begin{frame}{Lista Mínima}
  Cálculo Vol. 2 do Thomas 12\textsuperscript{a} ed. -- Seção 11.1
  \begin{enumerate}
    \item Estudar todo o texto da seção
    \item Resolver os exercícios: 3, 5, 11, 19, 21, 23, 31
  \end{enumerate}

  \vfill
  Atenção: A prova é baseada no livro, não nas apresentações
\end{frame}

%------------------------------------------------------------------------------%
\end{document}
%------------------------------------------------------------------------------%
